\documentclass[a4paper,12pt]{article}
\usepackage[utf8]{inputenc}
\usepackage{geometry}
\usepackage{listings}
\usepackage{hyperref}
\usepackage{booktabs}
\geometry{margin=2.5cm}

\title{Projekt 2: STRIPS}
\author{Paweł Adamczyk, Dominika Bujnarowska}
\date{}

\begin{document}

\maketitle

\section{Wstęp}

Celem projektu było zdefiniowanie dziedziny STRIPS oraz rozwiązanie problemów planowania z wykorzystaniem algorytmu forward planning. Dodatkowo zastosowano heurystykę oraz dekompozycję problemu na podcele w celu poprawy wydajności.

Rozwiązano trzy problemy o rosnącej złożoności (3, 4 i 5 miejsc), a każdy z nich analizowano dla czterech wariantów:
\begin{itemize}
    \item bez heurystyki,
    \item z heurystyką,
    \item z podcelami,
    \item z podcelami i heurystyką.
\end{itemize}

Kod źródłowy projektu dostępny jest w repozytorium GitHub:

\url{https://github.com/adpawel/gierki.git}

\section{Opis dziedziny}

Rozważana dziedzina to \textbf{birthday-dinner}.

\subsection{Własności}

Dla każdego miejsca $x$ zdefiniowano:
\begin{itemize}
    \item \texttt{clean(x)} -- miejsce jest czyste,
    \item \texttt{quiet(x)} -- miejsce jest spokojne,
    \item \texttt{garbage(x)} -- miejsce zawiera śmieci,
    \item \texttt{dinner(x)} -- kolacja przygotowana,
    \item \texttt{present(x)} -- prezent przygotowany.
\end{itemize}

\subsection{Akcje}

Zdefiniowano cztery akcje:

\begin{itemize}
    \item \textbf{cook(x)}  
    pre: \texttt{clean(x)}  
    eff: \texttt{dinner(x)}

    \item \textbf{wrap(x)}  
    pre: \texttt{quiet(x)}  
    eff: \texttt{present(x)}

    \item \textbf{carry(x)}  
    pre: \texttt{garbage(x)}  
    eff: $\neg$\texttt{garbage(x)}, $\neg$\texttt{clean(x)}

    \item \textbf{dolly(x)}  
    pre: \texttt{garbage(x)}  
    eff: $\neg$\texttt{garbage(x)}, $\neg$\texttt{quiet(x)}
\end{itemize}

\section{Heurystyka}

Zastosowana heurystyka:

\begin{quote}
liczba niespełnionych warunków celu
\end{quote}

\begin{verbatim}
h(s) = liczba warunków celu niespelnionych w stanie s
\end{verbatim}

Heurystyka przyspiesza wyszukiwanie, jednak w małych problemach jej wpływ jest ograniczony.

\section{Zadania na 4 i 6 punktów}

\subsection{Problemy}

Rozważono trzy problemy:
\begin{itemize}
    \item Problem 1: 3 miejsca (6 akcji),
    \item Problem 2: 4 miejsca (8 akcji),
    \item Problem 3: 5 miejsc (10 akcji).
\end{itemize}

\subsection{Osiągalne stany}

\begin{itemize}
    \item $4^3 = 64$,
    \item $4^4 = 256$,
    \item $4^5 = 1024$.
\end{itemize}

Wyniki zostały potwierdzone eksperymentalnie (BFS).

\subsection{Porównanie metod}

\begin{table}[h!]
\centering
\begin{tabular}{llccc}
\toprule
Problem & Metoda & Koszt & Czas [s] \\
\midrule
3 miejsca & bez heurystyki & 6 & 0.884 \\
3 miejsca & heurystyka & 6 & 0.868\\
3 miejsca & podcele & 6 & 1.711\\
3 miejsca & podcele + heurystyka & 6 & 1.689 \\
\midrule
4 miejsca & bez heurystyki & 8 & 0.897\\
4 miejsca & heurystyka & 8 & 0.869\\
4 miejsca & podcele & 8 & 2.572 \\
4 miejsca & podcele + heurystyka & 8 & 2.602 \\
\midrule
5 miejsc & bez heurystyki & 10 & 2.021\\
5 miejsc & heurystyka & 10 & 0.873\\
5 miejsc & podcele & 10 & 2.555 \\
5 miejsc & podcele + heurystyka & 10 & 2.562\\
\bottomrule
\end{tabular}
\caption{Porównanie wszystkich metod dla trzech problemów}
\end{table}

Podcele w tych problemach nie przyspieszyły działania – wprowadzały dodatkowy narzut obliczeniowy.

\subsection{Przykładowy plan (5 miejsc)}

\begin{verbatim}
1. cook_e
2. wrap_b
3. cook_c
4. wrap_a
5. wrap_c
6. wrap_e
7. cook_d
8. cook_a
9. wrap_d
10. cook_b
\end{verbatim}

\clearpage
\section{Zadania na 8 punktów}

Rozważono trudniejsze problemy o długości planu 20--24 akcje.

Wprowadzono ograniczenie czasu:
\begin{center}
\textbf{timeout = 500 s}
\end{center}

\subsection{Porównanie metod}

\begin{table}[h!]
\centering
\begin{tabular}{lcccc}
\toprule
Problem & Akcje & Metoda & Czas [s] & Wynik \\
\midrule
4 & 20 & Bez heurystyki & $>500$ & timeout \\
4 & 20 & Heurystyka & $>500$ & timeout \\
4 & 20 & Podcele & $>500$ & timeout \\
4 & 20 & Podcele + heurystyka & 3.59 & OK \\
\midrule
5 & 22 & Bez heurystyki & $>500$ & timeout \\
5 & 22 & Heurystyka & $>500$ & timeout \\
5 & 22 & Podcele & $>500$ & timeout \\
5 & 22 & Podcele + heurystyka & 3.01 & OK \\
\midrule
6 & 24 & Bez heurystyki & $>500$ & timeout \\
6 & 24 & Heurystyka & $>500$ & timeout \\
6 & 24 & Podcele & $>500$ & timeout \\
6 & 24 & Podcele + heurystyka & 4.07 & OK \\
\bottomrule
\end{tabular}
\caption{Porównanie metod dla problemów złożonych}
\end{table}

\subsection{Przykładowy plan (problem 6)}

\begin{verbatim}
  1. cook_a
  2. carry_a
  3. carry_h
  4. wrap_h
  5. dolly_b
  6. carry_g
  7. wrap_g
  8. cook_b
  9. carry_j
  10. wrap_j
  11. dolly_d
  12. carry_i
  13. wrap_i
  14. cook_d
  15. dolly_c
  16. cook_c
  17. dolly_f
  18. cook_f
  19. dolly_e
  20. cook_e
  21. carry_k
  22. wrap_k
  23. carry_l
  24. wrap_l
\end{verbatim}

\section{Wnioski}

\begin{itemize}
    \item Heurystyka poprawia wydajność dla większych problemów.
    \item W małych problemach jej wpływ jest niewielki.
    \item Podcele nie zawsze przyspieszają działanie – mogą wprowadzać narzut.
    \item W dużych problemach (≥20 akcji) konieczne jest połączenie heurystyki i podcelów.
    \item Metody bez heurystyki nie skalują się.
\end{itemize}

\section{Podsumowanie}

Projekt pokazuje, że efektywność planowania STRIPS zależy od rozmiaru problemu. Proste problemy można rozwiązać bez heurystyk, natomiast dla większych konieczne są dodatkowe techniki optymalizacji.

\end{document}